% Part: intuitionistic-logic
% Chapter: tableaux
% Section: soundness
%READERNOTE{SOL6-A207: मूळ adjacent-only अर्थनिर्धारणातील gap दुरुस्त करून सर्व वापरलेल्या पूर्वचिन्ह-विस्तारांसाठी प्राप्यता आवश्यक केली. नवा witness prefix घेताना सर्व जुने ancestors आणि जगांचे अर्थनिर्धारण जपले जाते हे सिद्ध केले; OLINC-233 ही ऐतिहासिक स्रोतदोषाची नोंद जपली.}

\documentclass[../../../include/open-logic-section]{subfiles}

\begin{document}

\olfileid{int}{tab}{sou}

\olsection{अंतःप्रज्ञावादी \usetoken{P}{tableau} ची निर्दोषता}

\begin{explain}
  अंतःप्रज्ञावादी !!{tableau}s निर्दोष आहेत हे
  दाखवण्यासाठी पुढील गृहीतकांना
  \[
  \sFmla{\True}{!B_1}[1], \dots, \sFmla{\True}{!B_n}[1], \sFmla{\False}{!A}[1]
  \]
  बंद !!{tableau} असेल, तर $!B_1, \dots, !B_n \Entails !A$
  हे दाखवावे लागेल. व्यत्यस्त रूप सिद्ध करणे सोपे आहे:
  एखाद्या $\mModel{M}$ प्रतिमानात आणि जगात~$w$,
  सर्व $i=1$, \dots,~$n$ साठी
  $\mSat{M}{!B_i}[w]$ असेल, पण
  $\mSat/{M}{!A}[w]$ असेल, तर कोणताही
  !!{tableau} बंद होऊ शकत नाही. असे
  प्रतिप्रतिमान !!{tableau} ची सुरुवातीची
  गृहीतके पूर्ततायोग्य असल्याचे दाखवते.
  पुराव्याची युक्ती अशी: !!a{tableau}
  च्या शाखेवरील सर्व पूर्वचिन्हयुक्त
  !!{formula}s पूर्ततायोग्य असतील, तर
  कोणताही नियम लावल्यानंतर मिळणाऱ्या
  विस्तारित शाखांपैकी किमान एक शाखा
  पूर्ततायोग्य असते. बंद शाखा अपूर्ततायोग्य
  असल्याने, पूर्ततायोग्य पूर्वचिन्हयुक्त
  !!{formula}s च्या संचावरील कोणत्याही
  !!{tableau} मध्ये किमान एक खुली शाखा
  असलीच पाहिजे.

  ही युक्ती अंतःप्रज्ञावादी प्रकरणात
  वापरण्यासाठी ‘पूर्ततायोग्य’ या संकल्पनेची
  व्याख्या प्राप्यता-संबंध आणि पूर्वचिन्हे
  यांना सामावून घेईल अशी वाढवावी लागते.
  त्यानंतर पुरावा सरळ आहे.
\end{explain}

\begin{defn}
  $P$ हा पूर्वचिन्हांचा एखादा संच असू द्या,
  म्हणजे $P \subseteq (\PosInt)^*
  \setminus \{\emptyseq\}$; तसेच
  $\mModel{M}$ हे प्रतिमान असू द्या.
  $f\colon P\to W$ हे फलन $\mModel{M}$ मधील
  $P$ चे \emph{अर्थनिर्धारण} असते तेव्हा आणि केवळ तेव्हाच,
  जेव्हा प्रत्येक $\sigma,\tau\in P$ साठी $\sigma$ हा
  $\tau$ चा आरंभीचा खंड असेल, तर $Rf(\sigma)f(\tau)$.
  ही अट सर्व विस्तारांसाठी आहे; मधली पूर्वचिन्हे
  $P$ मध्ये असणे आवश्यक नाही.

  पूर्वचिन्हे~$P$ यांच्या अर्थनिर्धारणाच्या
  संदर्भात पुढील व्याख्या करता येतात:
  \begin{enumerate}
  \item $\mModel{M}$ हे $\sFmla{\True}{!A}[\sigma]$
    पूर्ण करते तेव्हा आणि केवळ तेव्हाच
    $\mSat{M}{!A}[f(\sigma)]$.
  \item $\mModel{M}$ हे $\sFmla{\False}{!A}[\sigma]$
    पूर्ण करते तेव्हा आणि केवळ तेव्हाच
    $\mSat/{M}{!A}[f(\sigma)]$.
  \end{enumerate}
\end{defn}

%READERNOTE{OLINC-233: मूळ adjacent-prefix व्याख्येमुळे मधली पूर्वचिन्हे नसताना दीर्घ विस्तारापर्यंत प्राप्यता मिळत नव्हती. सध्याच्या मराठी व्याख्येत सर्व पूर्वचिन्ह-विस्तारांसाठी प्राप्यता थेट आवश्यक केली आहे; खालील बंदता-सिद्धतेला आता ती अट उपलब्ध आहे.}
या व्याख्येनुसार $\sigma$ आणि $\sigma.{*}$ ही दोन्ही
पूर्वचिन्हे वापरलेली असतील, तर $Rf(\sigma)f(\sigma.{*})$.
रिकाम्या विस्तारासाठी हे $R$ च्या परावर्तितेनेही मिळते.

\begin{defn}
  $\Gamma$ हा पूर्वचिन्हयुक्त !!{formula}s
  चा संच असू द्या आणि त्यात आलेल्या
  पूर्वचिन्हांचा संच $P(\Gamma)$ असू द्या.
  $f$ ही $P(\Gamma)$ ची
  $\mModel{M}$ मधील अर्थनिर्धारण
  असेल, तर $\mModel{M}$ हे
  $f$ च्या संदर्भात $\Gamma$
  पूर्ण करते, म्हणजे
  $\mSat{M}{\Gamma}[f]$, असे
  तेव्हा म्हणतो जेव्हा
  $\mModel{M}$ हे $\Gamma$ मधील
  प्रत्येक पूर्वचिन्हयुक्त
  !!{formula} हे~$f$ च्या संदर्भात
  पूर्ण करते. $\mModel{M}$ असे
  प्रतिमान आणि $P(\Gamma)$ ची
  $f$ असे अर्थनिर्धारण अस्तित्वात
  असतील की $\mSat{M}{\Gamma}[f]$,
  तेव्हा आणि केवळ तेव्हाच $\Gamma$ \emph{पूर्ततायोग्य}
  आहे.
\end{defn}

\begin{prop}
  $\Gamma$ मध्ये एखाद्या
  !!{formula}~$!A$ आणि पूर्वचिन्ह~$\sigma$
  साठी $\sFmla{\True}{!A}[\sigma]$ आणि
  $\sFmla{\False}{!A}[\sigma.{*}]$
  ही दोन्ही असतील, किंवा
  $\sFmla{\True}{\lfalse}[\sigma]$
  असेल, तर $\Gamma$ अपूर्ततायोग्य आहे.
\end{prop}

\begin{proof}
  नेहमीच $\mSat/{M}{\lfalse}[f(\sigma)]$
  असल्याने, $\sFmla{\True}{\lfalse}$
  असलेला $\Gamma$ अपूर्ततायोग्य असतो.

  आता दोन्ही चिन्हांकित सूत्रे पूर्ण
  करणारे $\mModel{M}$ प्रतिमान आणि
  $P(\Gamma)$ चे $f$ अर्थनिर्धारण
  असू शकत नाहीत. कारण
  $\mSat{M}{!A}[f(\sigma)]$
  असेल, तर
  \olref[sem][rel]{prop:true-monotonic}
  आणि $Rf(\sigma)f(\sigma.{*})$
  यांवरून
  $\mSat{M}{!A}[f(\sigma.{*})]$
  मिळते. म्हणून
  $\mSat{M}{!A}[f(\sigma)]$ आणि
  $\mSat/{M}{!A}[f(\sigma.{*})]$
  ही दोन्ही एकत्र असू शकत नाहीत.
\end{proof}

\begin{thm}[निर्दोषता]
  \ollabel{thm:tableau-soundness}
  $\Gamma$ ला बंद !!{tableau}
  असेल, तर $\Gamma$ अपूर्ततायोग्य
  आहे.
\end{thm}

\begin{proof}
!!a{tableau} च्या एखाद्या शाखेवरील
!!{signed formula}s चा संच
पूर्ततायोग्य असेल, तर ती शाखा
पूर्ततायोग्य म्हणू. किमान एक
पूर्ततायोग्य शाखा असलेला
!!a{tableau} पूर्ततायोग्य म्हणू.

आपण पुढील विधान दाखवू:
पूर्ततायोग्य !!{tableau} ला अनुमानाचा
कोणताही एक नियम लावल्यावरही
पूर्ततायोग्य !!{tableau} मिळतो.
यावरून प्रमेय सिद्ध होते: कोणताही
बंद !!{tableau} हा~$\Gamma$
मधील गृहीतके एवढीच असलेल्या
!!{tableau} ला नियम लावून मिळतो.
म्हणून $\Gamma$ पूर्ततायोग्य
असेल, तर त्यावरील प्रत्येक
!!{tableau} पूर्ततायोग्य असेल.
पण बंद !!{tableau} पूर्ततायोग्य
नसतो, कारण त्याच्या सर्व शाखा
बंद असतात आणि बंद शाखा
अपूर्ततायोग्य असतात.

आता पूर्ततायोग्य !!{tableau},
म्हणजे किमान एक पूर्ततायोग्य शाखा
असलेला !!a{tableau}, घेऊ.
नियम लावताना एखाद्या शाखेला
!!{signed formula}s जोडली जातात
किंवा तिचे दोन शाखांत विभाजन
होते. नियम न लावलेली एखादी
पूर्ततायोग्य शाखा आधीपासून
असेल, तर ती विस्तारित
!!{tableau} मध्येही तशीच
राहते आणि विस्तारित !!{tableau}
पूर्ततायोग्य असतो. म्हणून
पूर्ततायोग्य शाखेलाच नियम
लावल्याचे प्रकरण पाहणे पुरेसे आहे.

त्या शाखेवरील !!{signed formula}s
चा संच $\Gamma$ असू द्या आणि
नियम ज्या !!{signed formula}
ला लावतो ते
$\sFmla{S}{!A}[\sigma] \in \Gamma$
असू द्या. नियमामुळे शाखा फुटत
नसेल, तर नियमाच्या निष्कर्षांसह
$\Gamma$ असलेली विस्तारित शाखा
अजूनही पूर्ततायोग्य आहे हे
दाखवावे लागते. शाखा फुटत
असेल, तर मिळणाऱ्या दोन शाखांपैकी
किमान एक पूर्ततायोग्य आहे हे
दाखवावे लागते.
\begin{enumerate}
\item $\sFmla{\True}{\lnot !B}[\sigma] \in \Gamma$
  याला $\TRule{\True}{\lnot}$ लावून
  शाखा वाढवू. मग विस्तारित
  शाखेत $\Gamma \cup
  \{\sFmla{\False}{!B}[\sigma.{*}]\}$
  ही !!{signed formula}s असतात.
  $\mSat{M}{\Gamma}[f]$ असे समजा.
  विशेषतः
  $\mSat{M}{\lnot !B}[f(\sigma)]$.
  म्हणून $Rf(\sigma)w$ असलेल्या
  प्रत्येक $w$ साठी
  $\mSat/{M}{!B}[w]$; यात
  $f(\sigma.{*})$ सुद्धा येते.
  त्यामुळे $\mModel{M}$ हे
  $\sFmla{\False}{!B}[\sigma.{*}]$
  $f$ च्या संदर्भात पूर्ण करते.
\item $\sFmla{\False}{\lnot !B}[\sigma] \in \Gamma$
  याला $\TRule{\False}{\lnot}$
  लावून शाखा वाढवण्याचे प्रकरण:
  सरावासाठी.
\item $\sFmla{\True}{!B \land !C}[\sigma] \in \Gamma$
  याला $\TRule{\True}{\land}$
  लावल्यावर त्याच शाखेत
  $\sFmla{\True}{!B}[\sigma]$
  आणि $\sFmla{\True}{!C}[\sigma]$
  ही दोन नवी !!{signed formula}s
  मिळतात. $\mSat{M}{\Gamma}[f]$
  असे समजा; विशेषतः
  $\mSat{M}{!B \land
    !C}[f(\sigma)]$. मग
  $\mSat{M}{!B}[f(\sigma)]$
  आणि $\mSat{M}{!C}[f(\sigma)]$.
  म्हणजे $\mModel{M}$ हे
  $\sFmla{\True}{!B}[\sigma]$
  आणि $\sFmla{\True}{!C}[\sigma]$
  दोन्ही~$f$ च्या संदर्भात
  पूर्ण करते.
\item $\sFmla{\False}{!B \lor !C}[\sigma] \in \Gamma$
  याला $\TRule{\False}{\lor}$
  लावण्याचे प्रकरण: सरावासाठी.
\item $\sFmla{\False}{!B \lif !C}[\sigma] \in \Gamma$
  याला $\TRule{\False}{\lif}$
  लावल्यावर शाखेत
  $\sFmla{\True}{!B}[\sigma.n]$
  आणि $\sFmla{\False}{!C}[\sigma.n]$
  ही दोन नवी !!{signed formula}s
  मिळतात. येथे $\sigma.n$
  हे शाखेसाठी नवे पूर्वचिन्ह आहे: $\sigma.n\notin P(\Gamma)$,
  आणि त्याचा कोणताही विस्तारही $P(\Gamma)$ मध्ये नाही.

  $\Gamma$ पूर्ततायोग्य असल्याने,
  $\mSat{M}{\Gamma}[f]$ असेल
  असे $\mModel{M}$ प्रतिमान
  आणि $P(\Gamma)$ ची~$f$
  अर्थनिर्धारण आहेत. विशेषतः
  $\mSat/{M}{!B \lif !C}[f(\sigma)]$.
  पुढील संच पूर्ततायोग्य आहे हे
  दाखवावे लागेल:
  $\Gamma \cup
  \{\sFmla{\True}{!B}[\sigma.n],
  \sFmla{\False}{!C}[\sigma.n]\}$.
  त्यासाठी $P(\Gamma)
  \cup \{\sigma.n\}$
  याचे अर्थनिर्धारण पुढीलप्रमाणे
  ठरवू:

  $\mSat/{M}{!B \lif !C}[f(\sigma)]$
  असल्याने $Rf(\sigma)w$,
  $\mSat{M}{!B}[w]$ आणि
  $\mSat/{M}{!C}[w]$
  असा $w \in W$ असतो.
  $f'$ हे $f$ सारखेच, फक्त
  $f'(\sigma.n) = w$
  असे ठरवू. $f'(\sigma) = f(\sigma)$
  आणि $Rf(\sigma)w$ असल्याने $Rf'(\sigma)f'(\sigma.n)$.
  जुन्या संचातील $\tau$ हा $\sigma.n$ चा आरंभीचा खंड
  असेल, तर तो $\sigma$ चाही आरंभीचा खंड आहे.
  म्हणून जुन्या अर्थनिर्धारणाने $Rf(\tau)f(\sigma)$,
  आणि संक्रमणामुळे $Rf'(\tau)f'(\sigma.n)$.
  नव्या पूर्वचिन्हाचा कोणताही विस्तार जुन्या संचात नाही;
  त्यामुळे त्यापासून जुन्या विस्ताराकडे अतिरिक्त अट येत नाही.
  नव्या पूर्वचिन्हाची स्वतःशी अट परावर्तितेने मिळते.
  म्हणून $f'$ हे $P(\Gamma)\cup\{\sigma.n\}$ चे
  सर्व पूर्वचिन्ह-विस्तारांसाठी अर्थनिर्धारण आहे.
  तसेच $\mSat{M}{!B}[f'(\sigma.n)]$
  आणि $\mSat/{M}{!C}[f'(\sigma.n)]$.
  सर्व पूर्वचिन्हे
  $\sigma' \in P(\Gamma)$ साठी
  $f(\sigma') = f'(\sigma')$
  असल्याने
  $\mSat{M}{\Gamma}[f']$.
  म्हणून $\mModel{M}, f'$ हे
  $\Gamma \cup
  \{\sFmla{\True}{!B}[\sigma.n],
  \sFmla{\False}{!C}[\sigma.n]\}$
  पूर्ण करते.
\end{enumerate}
आता दोन शाखा निर्माण करणारे
नियम पाहू.
\begin{enumerate}
\item $\sFmla{\False}{!B \land !C}[\sigma] \in \Gamma$
  याला $\TRule{\False}{\land}$
  लावल्यावर दोन शाखा मिळतात:
  डावीकडील शाखेत
  $\sFmla{\False}{!B}[\sigma]$
  आणि उजवीकडील शाखेत
  $\sFmla{\False}{!C}[\sigma]$.
  $\mSat{M}{\Gamma}[f]$
  असे समजा; विशेषतः
  $\mSat/{M}{!B \land !C}[f(\sigma)]$.
  मग $\mSat/{M}{!B}[f(\sigma)]$
  किंवा $\mSat/{M}{!C}[f(\sigma)]$.
  पहिल्या प्रकरणात
  $\mModel{M}, f$ हे
  $\sFmla{\False}{!B}[\sigma]$
  पूर्ण करते; म्हणजे डावी
  शाखा पूर्ततायोग्य आहे.
  दुसऱ्या प्रकरणात
  $\mModel{M}, f$ हे
  $\sFmla{\False}{!C}[\sigma]$
  पूर्ण करते; म्हणजे उजवी
  शाखा पूर्ततायोग्य आहे.
\item $\sFmla{\True}{!B \lor !C}[\sigma] \in \Gamma$
  याला $\TRule{\True}{\lor}$
  लावण्याचे प्रकरण: सरावासाठी.
\item $\sFmla{\True}{!B \lif !C}[\sigma] \in \Gamma$
  याला $\TRule{\True}{\lif}$
  लावण्याचे प्रकरण: सरावासाठी.
\end{enumerate}
\end{proof}

\begin{prob}
\olref[int][tab][sou]{thm:tableau-soundness}
या प्रमेयाचा पुरावा पूर्ण करा.
\end{prob}

\begin{cor}
\ollabel{cor:entailment-soundness}
$\Gamma \Proves !A$ असेल,
तर $\Gamma \Entails !A$.
\end{cor}

\begin{proof}
  $\Gamma \Proves !A$ असेल,
  तर काही
  $!B_1$, \dots, $!B_n \in
  \Gamma$ साठी
  $\Delta = \{\sFmla{\False}{!A}[1], \sFmla{\True}{!B_1}[1],
  \dots, \sFmla{\True}{!B_n}[1]\}$
  याला बंद !!{tableau} असतो.
  $\Gamma \Entails !A$
  हे दाखवायचे आहे. तसे नाही
  असे समजा. मग एखाद्या
  $\mModel{M}$ प्रतिमानात
  आणि जगात~$w$,
  $i=1$, \dots,~$n$ साठी
  $\mSat{M}{!B_i}[w]$,
  पण $\mSat/{M}{!A}[w]$.
  $f(1) = w$ असे ठरवू;
  मग $f$ ही
  $P(\Delta)$ ची
  $\mModel{M}$ मधील
  अर्थनिर्धारण आहे आणि
  $\mModel{M}$ हे
  $\Delta$ ला~$f$ च्या
  संदर्भात पूर्ण करते.
  पण \olref{thm:tableau-soundness}
  नुसार, बंद !!{tableau}
  असल्याने $\Delta$
  अपूर्ततायोग्य आहे; ही
  विसंगती आहे. म्हणून
  $\Gamma \Entails !A$
  हेच मिळते.
\end{proof}

\begin{cor}
\ollabel{cor:weak-soundness}
$\Proves !A$ असेल, तर
$!A$ हे सर्व प्रतिमानांत
सत्य असते.
\end{cor}

\end{document}
